% ============================================================
%  Chapter 13 — Null Space and Kernel
% ============================================================
\section{Null Space and Kernel}

Given a matrix $A$, the set of all vectors that $A$ maps to zero is a
subspace called the \emph{null space} or \emph{kernel} of $A$.  Its dimension
(the \emph{nullity}) together with the rank of $A$ adds up to the number of
columns — the content of the Rank-Nullity Theorem.

% ---- Definition --------------------------------------------
\begin{defbox}[Definition --- Null Space (Kernel) of a Matrix]
\label{def:null-space-matrix}
For $A\in\Mmn{m}{n}$, the \textbf{null space} (or \textbf{kernel}) of $A$ is
\[
    \Ker(A) = \{v\in\F^n \mid Av = \mathbf{0}_m\}.
\]
Its dimension is the \textbf{nullity} of $A$:
$\operatorname{nullity}(A) = \dime_\F\Ker(A)$.
\end{defbox}

\begin{tipbox}[Remark --- Kernel is a Subspace of $\F^n$]
$\Ker(A)$ is the solution set of the homogeneous linear system $Ax=\mathbf{0}$,
which we showed in Chapter~7 (Example: solution set of $Ax=\mathbf{0}$) is
always a subspace of $\F^n$.
\end{tipbox}

% ---- Finding a basis for the kernel via RREF ---------------
\begin{methodbox}[Method --- Finding a Basis for the Null Space]
\label{meth:null-space-basis}
Given $A\in\Mmn{m}{n}$:
\begin{enumerate}
    \item Row-reduce $A$ to its RREF $R$.
    \item Identify the \emph{free variables} (non-pivot columns
          $j_1,\ldots,j_{n-r}$ where $r=\rk(A)$).
    \item For each free variable, set it to $1$ and all other free variables
          to $0$; solve for the pivot variables.  The resulting vector is a
          basis vector of $\Ker(A)$.
    \item $\Ker(A) = \spn\{\text{these } n-r \text{ vectors}\}$ and they are l.i.
\end{enumerate}
\end{methodbox}

% ---- Example 1 ---------------------------------------------
\begin{exbox}[Example --- Kernel of a $4\times 4$ Matrix (RREF Given)]
Let
\[
    A = \begin{pmatrix}
        1 & 0 & 0 & -1 \\
        0 & 1 & 0 &  3 \\
        0 & 0 & 1 &  2 \\
        0 & 0 & 0 &  0
    \end{pmatrix}
    \quad(\text{already in RREF}).
\]
A vector $(x_1,x_2,x_3,x_4)\in\Ker(A)$ satisfies
\[
    x_1 = x_4,\quad x_2 = -3x_4,\quad x_3 = -2x_4,
\]
with $x_4$ free.  Setting $x_4=1$:
\[
    \Ker(A) = \spn\!\left\{(1,-3,-2,1)\right\},
    \qquad \operatorname{nullity}(A)=1.
\]
\end{exbox}

% ---- Example 2 ---------------------------------------------
\begin{exbox}[Example --- Kernel of a $3\times 4$ Matrix (RREF Given)]
Let
\[
    A = \begin{pmatrix}
        1 & 3 & 0 & 1 \\
        0 & 0 & 1 & 2 \\
        0 & 0 & 0 & 0
    \end{pmatrix}
    \quad(\text{RREF; pivots in columns 1 and 3}).
\]
Free variables: $x_2$ and $x_4$.  The constraints are
$x_1 = -3x_2 - x_4$ and $x_3 = -2x_4$.  Setting $(x_2,x_4)=(1,0)$ and
$(0,1)$ in turn:
\[
    \Ker(A) = \spn\!\left\{(-3,1,0,0),\,(-1,0,-2,1)\right\},
    \qquad \operatorname{nullity}(A)=2.
\]
\end{exbox}

% ---- Rank-Nullity Theorem (matrix version) -----------------
\begin{thmbox}[Theorem --- Rank-Nullity Theorem (Matrix Version)]
\label{thm:rank-nullity-matrix}
For any $A\in\Mmn{m}{n}$:
\[
    \rk(A) + \dime_\F\Ker(A) = n.
\]
That is, the rank plus the nullity of $A$ equals the number of columns.
\end{thmbox}

\begin{proofbox}[Proof --- Rank-Nullity Theorem (Matrix Version)]
Let $R$ be the RREF of $A$ and let $r=\rk(A)$ be the number of pivot columns.
Then $R$ (and hence $A$) has $n-r$ non-pivot (free) columns.  The method above
produces exactly $n-r$ linearly independent vectors that span $\Ker(A)$, so
$\dime\Ker(A)=n-r$.  Therefore $\rk(A)+\dime\Ker(A)=r+(n-r)=n$.
\hfill$\blacksquare$
\end{proofbox}

% ---- Summary box -------------------------------------------
\begin{flowbox}[Key Formula --- Rank and Nullity]
\[
    \boxed{\rk(A) + \operatorname{nullity}(A) = n}
    \qquad \text{for } A\in\Mmn{m}{n}.
\]
\begin{itemize}
    \item $\rk(A)$ = number of pivot columns = $\dime(\text{column space of }A)$.
    \item $\operatorname{nullity}(A)$ = number of free variables
          = $\dime\Ker(A)$.
\end{itemize}
\end{flowbox}
